\section{Giới thiệu}\label{sec:intro}

Các môn lập trình giao rất nhiều bài tập nhỏ, và người học tiến bộ với tốc độ rất khác nhau. Một bài quá khó làm người học nản lòng; một bài quá dễ dạy được rất ít. Gợi ý bài tập thích ứng giải quyết vấn đề này bằng cách ước lượng trạng thái của từng người học từ các lần làm trước và chọn bài tiếp theo cho phù hợp — một ý tưởng thường được gắn với vùng phát triển gần \citep{vygotsky1978mind} và có truyền thống lâu dài trong các hệ thống giáo dục thích ứng cho lập trình \citep{brusilovsky2007user,hsiao2010guiding}. Một công thức nay đã phổ biến kết hợp mô hình truy vết tri thức (knowledge tracing, KT), vốn dự báo khả năng người học làm đúng mỗi bài từ lịch sử của họ, với một tác tử học tăng cường (reinforcement learning, RL) coi việc chọn bài là một bài toán quyết định tuần tự: trạng thái của mô hình người học là trạng thái của tác tử, gợi ý một bài là một hành động, và mức tăng tri thức mà mô hình dự đoán là phần thưởng \citep{huang2019dre,liu2019cseal,ai2019cdkt,chen2023src,li2023gehrl,ozyurt2025exrec}. Vì thử nghiệm với người học tốn kém, các tác tử như vậy thường được huấn luyện \emph{và đánh giá} ngay bên trong mô hình người học đã fit, và mức tăng tri thức dự đoán được báo cáo như lợi ích của chính sách \citep{cheng2026unier}.

Cách đánh giá này chỉ tốt bằng mức bộ mô phỏng mô tả đúng việc luyện tập làm thay đổi khả năng làm đúng các bài khác như thế nào. Mô hình KT được fit và kiểm định trên dự báo bước kế tiếp của những chuỗi mà người học đã thực sự đi qua. Một chính sách, ngược lại, tạo ra chuỗi mới, và bộ mô phỏng phải ngoại suy hệ quả của các chuỗi đó. Hai mô hình có thể dự báo dữ liệu quan sát tốt như nhau mà vẫn bất đồng về những phản thực tế (counterfactual) như vậy \citep{breiman2001two,damour2022underspecification}. Với các quyết định dạy học, điều này đã được nhận ra từ sớm: tham số truy vết tri thức cá nhân hóa và tham số chung cho cả quần thể hàm ý lượng luyện tập rất khác nhau \citep{lee2012individualizing}, và các mô hình người học có độ chính xác dự báo tương tự có thể đề xuất lượng luyện tập khác nhau đáng kể \citep{rollinson2015predictive}. Ma trận đánh giá bền vững vì thế đánh giá mỗi chính sách ứng viên dưới nhiều mô hình người học \citep{doroudi2017rem}, và một tổng quan về RL cho sắp xếp trình tự dạy học nhận thấy các chính sách rút ra từ mô hình thuần dữ liệu thường không vượt được các đường cơ sở hợp lý khi thử với học sinh \citep{doroudi2019reward}. Trong học tăng cường nói chung, một bộ tối ưu được huấn luyện trên một mô hình hay một phần thưởng học từ dữ liệu sẽ khai thác chính sai số của mô hình đó \citep{kurutach2018metrpo,janner2019trust,gao2023scaling}. Các pipeline KT--RL hiện đại gộp cả hai rủi ro — mô hình học từ dữ liệu đồng thời là môi trường, phần thưởng và thước đo — nhưng hiếm khi có ai định lượng xem bao nhiêu phần của tiến bộ được báo cáo phản ánh việc học mà bộ mô phỏng mã hóa, và bao nhiêu phản ánh các đặc điểm riêng của bộ mô phỏng.

Chúng tôi nghiên cứu câu hỏi này cho lập trình nhập môn, dùng {{n_submissions}} lượt nộp C/C++ đã được chấm của {{n_learners}} người học trên hệ thống chấm bài trực tuyến ACcoding \citep{chen2024accoding}. Chúng tôi fit sáu mô hình người học có độ chính xác dự báo ngang nhau, gồm DKT (truy vết tri thức sâu) làm tracker triển khai và một cây định danh dựa trên entropy để chẩn đoán dễ hiểu; huấn luyện các bộ gợi ý Double DQN và PPO trong ba loại thế giới huấn luyện; và đánh giá mọi chính sách trong chín thế giới khác nhau về tập người học được dùng để fit và về cách biểu diễn việc học. Hai công cụ làm cho việc đánh giá có tính chẩn đoán. Một thế giới \emph{đối chứng âm} không có bất kỳ việc học nào, trong đó tiến bộ thật của mọi chính sách đúng bằng 0, phơi bày trực tiếp các tiến bộ giả. Một \emph{phân rã} chính xác tiến bộ mà tracker báo cáo thành thay đổi dự báo mà mô hình tự kỳ vọng (\emph{độ trôi}) và \emph{phần bất ngờ} của bằng chứng cho thấy tiến bộ đó đến từ đâu. Chúng tôi chứng minh rằng độ trôi triệt tiêu với mô hình Bayes nhất quán không có học, và rằng bên trong chính bộ mô phỏng của nó, tiến bộ kỳ vọng của mọi chính sách bằng độ trôi mà chính sách tích lũy; do đó đánh giá trong mô phỏng thưởng cho chính sách vì đã chọn những bài mà mô hình \emph{kỳ vọng} sẽ nâng dự báo của chính nó.

Chúng tôi trả lời bốn câu hỏi nghiên cứu. RQ1: Các mô hình trạng thái nhận thức dự báo khả năng làm đúng ở lần nộp đầu của người học lập trình chưa thấy tốt đến đâu, và chúng có khác nhau không? RQ2: Tiến bộ của các bộ gợi ý RL lớn đến đâu trong bộ mô phỏng KT mà chúng được huấn luyện? RQ3: Các tiến bộ đó có đứng vững trong các thế giới fit độc lập không, và tracker báo tiến bộ bao nhiêu khi việc học là không thể? RQ4: Điều gì tạo ra sự chênh lệch, và chính sách nào vẫn bền vững? Tất cả protocol, gồm bốn giả thuyết, đã được khóa trên một cohort phát triển trước khi nạp quần thể đánh giá.

Đóng góp của bài báo gồm: (1) một pipeline đầu-cuối, có thể tái lập, cho gợi ý bài tập lập trình thích ứng trên dữ liệu thật, từ mô hình người học đến các tác tử DQN và PPO đã huấn luyện và một ứng dụng tương tác, với các giả thuyết đăng ký trước được kiểm định trên một quần thể chưa từng dùng; (2) bằng chứng rằng dù sáu mô hình người học dự báo tốt ngang nhau, tiến bộ trong mô phỏng của các tác tử thưởng theo KT co lại gần 0 trong các thế giới fit độc lập có học ẩn, và tương quan thứ hạng gần như hoàn hảo với tiến bộ mà tracker báo cáo trong một thế giới không có học; (3) phép phân rã độ trôi cùng hai mệnh đề giải thích các tiến bộ đó là thay đổi tự dự đoán tích lũy của tracker chứ không phải việc học — một chẩn đoán nội tại của mô hình, không cần thêm dữ liệu; và (4) bằng chứng thăm dò rằng tác tử huấn luyện trong thế giới có cấu trúc học tường minh chuyển giao tốt nhất giữa các thế giới ẩn cùng họ — nhưng không sang các thế giới nơ-ron — cùng các kiểm tra ít tốn kém cho đánh giá dựa trên mô phỏng.

\Cref{sec:related} điểm lại các nghiên cứu liên quan, \cref{sec:data} mô tả dữ liệu và \cref{sec:methods} trình bày các mô hình, bài toán quyết định, các thế giới đánh giá và phương pháp phân tích. \Cref{sec:results} báo cáo kết quả, \cref{sec:discussion} thảo luận hàm ý và giới hạn, và \cref{sec:conclusion} kết luận.
